\answer{ex:08:1}{(a)~$\approx1.7\cdot10^6$。(b)~$g=\frac{0.9}{\sqrt{0.19}}\,\sigma_{\text{後}}/\sigma_{\text{前}}\approx16.5$。答えは雑音の比だけで決まる。}
\answer{ex:08:2}{(a)~固有値は$-(1\pm g\beta)/\tau$。差は$\tau/(1-g\beta)$で減衰する。$g\beta=0.5$で$40\ms$（入力の差は3倍に増幅）、$0.9$で$200\ms$（$19$倍）。(b)~$0.905$と$1.105$。二つの境界の間では強い入力だけが勝つ。}
\answer{ex:08:3}{$P(k)=e^{gu_k/\sigma}/\sum_le^{gu_l/\sigma}$、つまり\eqref{eq:softmax}。$g=10\ln9\approx22$。}
\answer{ex:08:4}{$\bar\Delta=1/3$、$g^*\approx3\ln(1/h)$：$21$、$14$、$9$。モデルでは$16$--$23$、$11$、$8$。}
\answer{ex:08:5}{$g^*$は$h=0.001$の$16$--$23$から$h=0.2$の$6$--$8$まで。見積もり$3\ln(1/h)$は$h\le0.05$で1.5倍以内で正しい。$h\gtrsim\eta/10$では規則$Q\leftarrow Q+\eta(R-Q)$が探索で得たものを取り込む時間がなく、$g^*$は約$8$で止まる。}
\answer{ex:08:6}{$g_{\text{lo}}=0$で$T_p=\tau\ln9=44\ms$（モデルも$44\ms$）、$g_{\text{lo}}=0.25$で$70\ms$。ゲインをほぼ0にする$2$--$3\tau$のバーストで足りる。}
\answer{ex:08:7}{(a)~最良の一定値$0.925$（$g=8$）、オラクル$0.929$（$16/8$）。得られるものはほとんどない。(b)~たとえば報酬がまれな穏やかな時期（$p\in[0,0.3]$）と$h=0.1$の変化の多い時期。$0.850$対$0.872$。調整は$\eta=0.2$で約4分の1、$\eta=0.05$でほぼすべてを取り戻す。}
